\chapter{Θεωρητικό Υπόβαθρο και Σχετική Βιβλιογραφία}
\label{ch:background}

\section{Από το 5G NTN στην ενοποίηση TN-NTN στο 6G}
\label{sec:bg-3gpp-evolution}

Η ενσωμάτωση μη-επίγειων δικτύων στα πρότυπα κινητών επικοινωνιών ξεκίνησε
επίσημα με μελέτες σκοπιμότητας του 3GPP στα Release~15 και Release~16, και
έλαβε πλήρεις προδιαγραφές στο Release~17 (TR~38.821~\cite{tr38821}), το
οποίο ορίζει πρόσβαση 5G NR μέσω δορυφόρου σε ζώνες συχνοτήτων S και L, με
έμφαση σε \emph{transparent payload} αρχιτεκτονική (ο δορυφόρος λειτουργεί
ως αναμεταδότης ραδιοσυχνοτήτων, ενώ η λειτουργικότητα gNB παραμένει στο
έδαφος) και σε τερματικά χειρός.

Είναι σημαντικό να τονιστεί ότι, όπως περιγράφουν οι
Noh~κ.ά.~\cite{noh2024tnntn}, οι προδιαγραφές NTN μέχρι και το Release~19
προϋποθέτουν ότι ένα τερματικό χρήστη (UE) \emph{δεν} διατηρεί ταυτόχρονα
ζεύξη TN και NTN: η μία ζεύξη πρέπει να απενεργοποιηθεί πριν ολοκληρωθεί η
διαδικασία επανεπιλογής κυψέλης ή handover προς την άλλη. Δηλαδή, μηχανισμοί
όπως το carrier aggregation (CA) και το dual connectivity (DC), που είναι
πλήρως τυποποιημένοι για επίγεια δίκτυα, \emph{δεν} υποστηρίζονται ακόμα για
NTN. Η ταυτόχρονη ή δυναμικά εναλλασσόμενη λειτουργία TN και NTN --
αναφερόμενη ως \emph{TN-NTN aggregation} και \emph{TN-NTN unification}
αντίστοιχα -- παραμένει, επομένως, ανοιχτό ζήτημα που αναμένεται να
αντιμετωπιστεί σε μελλοντικά releases προς το 6G.

Ένα πρακτικό εμπόδιο για την ταυτόχρονη λειτουργία TN-NTN που αναδεικνύουν
οι Noh~κ.ά.~\cite{noh2024tnntn} είναι το μεγάλο εύρος ισχύος λήψης
(\emph{dynamic range}) που πρέπει να διαχειριστεί ο δέκτης ενός UE όταν
λαμβάνει ταυτόχρονα από επίγειο σταθμό βάσης και δορυφόρο: λόγω των πολύ
διαφορετικών επιπέδων EIRP και απωλειών διάδοσης μεταξύ TN, LEO σε
διαφορετικά υψόμετρα, και GEO, το εύρος αυτό μπορεί να φτάσει τα 70~dB στη
ζώνη S, απαιτώντας δέκτες υψηλής δυναμικής περιοχής ή μηχανισμούς
εξισορρόπησης ισχύος (power control) και στις δύο κατευθύνσεις. Το εύρημα
αυτό είναι ενδεικτικό του γενικότερου προβλήματος: η ετερογένεια μεταξύ
επίγειων και δορυφορικών ζεύξεων δεν είναι απλώς ζήτημα κάλυψης, αλλά
επηρεάζει και τον σχεδιασμό του hardware και των πρωτοκόλλων.

\section{Πολυσυνδεσιμότητα σε Ολοκληρωμένα Δορυφορικά-Επίγεια Δίκτυα}
\label{sec:bg-mc-stin}

\subsection{Αρχιτεκτονικές ανάπτυξης}
\label{subsec:bg-architectures}

Οι Li και Shang~\cite{li2025advancing} προτείνουν τρεις θεμελιώδεις
αρχιτεκτονικές ανάπτυξης πολυσυνδεσιμότητας σε STINs:

\begin{description}
  \item[Multi-Satellite (MS)] Πολλαπλοί δορυφόροι συντονίζονται από μια
    κεντρική μονάδα επεξεργασίας (CPU/CCS) εγκατεστημένη σε έναν από
    αυτούς, και μεταδίδουν από κοινού προς τον χρήστη στο ίδιο
    time-frequency resource block. Προσφέρει σχεδόν καθολική κάλυψη σε
    περιοχές χωρίς επίγεια υποδομή (ωκεανοί, πολικές περιοχές, καταστάσεις
    έκτακτης ανάγκης), με το κόστος αυξημένου λανθάνοντα χρόνου και
    πολυπλοκότητας διαχείρισης φάσματος/ισχύος μεταξύ δορυφόρων.
  \item[Single-Satellite Single-Base-Station (SS-SBS)] Ένας δορυφόρος και
    ένας επίγειος σταθμός βάσης είναι διαθέσιμοι για έναν χρήστη, με
    τέσσερις πιθανές καταστάσεις σύνδεσης (εξυπηρέτηση από τους δύο
    ταυτόχρονα, μόνο από τον δορυφόρο, μόνο από τον σταθμό βάσης, ή από
    κανέναν). Είναι η απλούστερη αρχιτεκτονική σε όρους συντονισμού και
    λανθάνοντα χρόνου, και αποτελεί το φυσικό σημείο εκκίνησης για
    πρώιμες αναπτύξεις STIN. \textbf{Το σενάριο της παρούσας εργασίας
    (δύο επίγειοι σταθμοί βάσης, ένας δορυφόρος LEO, single-connectivity
    ανά χρήστη) αποτελεί μια απλοποιημένη, πολλαπλών-BS παραλλαγή αυτής
    ακριβώς της αρχιτεκτονικής.}
  \item[Multi-Satellite Multi-Base-Station (MS-MBS)] Γενίκευση των δύο
    παραπάνω, με από κοινού συντονισμό της κεντρικής μονάδας NTN και της
    κεντρικής μονάδας TN. Προσφέρει τη μέγιστη χωρητικότητα και
    αξιοπιστία, αλλά και τη μεγαλύτερη πολυπλοκότητα υλοποίησης
    (συγχρονισμός, κατανομή πόρων, διαχείριση handover) και κόστος
    υποδομής.
\end{description}

\subsection{Προκλήσεις σχεδίασης σε επίπεδο συστήματος}
\label{subsec:bg-challenges}

Πέρα από την αρχιτεκτονική ανάπτυξης, οι~\cite{li2025advancing} εντοπίζουν
τέσσερις βασικές τεχνικές προκλήσεις για την πολυσυνδεσιμότητα σε STINs:

\begin{itemize}
  \item \textbf{Δικτύωση δορυφόρων} (cell-free ή multi-tier σχήματα
    συνεργασίας μεταξύ δορυφόρων).
  \item \textbf{Beamforming}, όπου η ταυτόχρονη κατεύθυνση δεσμών προς
    πολλαπλούς χρήστες από επίγειους και δορυφορικούς κόμβους περιπλέκεται
    από την κινητικότητα του δορυφόρου και τη μεταβλητή ζήτηση των
    χρηστών.
  \item \textbf{Εκτίμηση καναλιού}, όπου η μεγάλη και μεταβλητή καθυστέρηση
    διάδοσης δυσχεραίνει την απόκτηση επίκαιρης πληροφορίας κατάστασης
    καναλιού (CSI) και για τις δύο ζεύξεις.
  \item \textbf{Συγχρονισμός}, όπου οι καθυστερήσεις διάδοσης δορυφόρου και
    επίγειου σταθμού βάσης πρέπει να αντισταθμιστούν ώστε τα σήματα να
    φτάνουν στον χρήστη στο ίδιο χρονικό παράθυρο (π.χ. εντός του κυκλικού
    προθέματος).
\end{itemize}

Ο Majamaa~\cite{majamaa2024toward} εστιάζει ειδικά στην πολυσυνδεσιμότητα
επιπέδου PDCP (Multi-Radio Dual Connectivity - MR-DC), όπου ένα UE
συνδέεται σε έναν κόμβο-master (MN) και έναν κόμβο-secondary (SN) μέσω της
διεπαφής Xn. Ενώ το MR-DC είναι πλήρως τυποποιημένο για TN, η υποστήριξή
του για NTN δεν έχει ακόμα οριστικοποιηθεί (είχε προγραμματιστεί για το
Release~19 αλλά αφαιρέθηκε από το πεδίο εφαρμογής). Ο συγγραφέας συνοψίζει
επτά κατηγορίες προβλημάτων ειδικών για NTN MC, οι σημαντικότερες εκ των
οποίων για την παρούσα εργασία είναι:

\begin{itemize}
  \item \textbf{Κινητικότητα}: Η υψηλή ταχύτητα των δορυφόρων LEO προκαλεί
    συχνές αλλαγές/προσθήκες/αφαιρέσεις δευτερεύοντος κόμβου (SN), σε
    αντίθεση με τα στατικά επίγεια δίκτυα όπου η κινητικότητα προέρχεται
    μόνο από τον χρήστη. Προτεινόμενη λύση: προγνωστικά σχήματα βασισμένα
    σε δεδομένα εφημερίδας τροχιάς (ephemeris) του δορυφόρου και σε
    δεδομένα GNSS του χρήστη.
  \item \textbf{Κριτήριο προσθήκης δευτερεύοντος κόμβου}: Στα TN,
    χρησιμοποιείται συνήθως η ισχύς λαμβανόμενου σήματος αναφοράς. Στα
    NTN, λόγω επαναχρησιμοποίησης συχνότητας και μικρών διαφορών ισχύος
    μεταξύ γειτονικών δεσμών, προτείνονται εναλλακτικά κριτήρια όπως η
    θέση ή η γωνία ανύψωσης -- ακριβώς το κριτήριο (γωνία ανύψωσης άνω
    ενός ελάχιστου ορίου) που υιοθετεί και η προσομοίωση της παρούσας
    εργασίας (βλ.~Ενότητα~\ref{sec:meth-sat-channel}).
  \item \textbf{Κατανάλωση ισχύος}: Προτείνεται η χρήση της καλύτερης
    ζεύξης ανά πάσα στιγμή (best-link selection) στην ανοδική ζεύξη ως
    μηχανισμός μείωσης κατανάλωσης στο UE -- η ίδια λογική με τον greedy
    SNR κανόνα επιλογής κόμβου που υλοποιεί η παρούσα προσομοίωση, αν και
    εδώ εφαρμόζεται στην καθοδική ζεύξη και συνδυάζεται με ρητή
    μοντελοποίηση ενέργειας ανά bit.
\end{itemize}

\section{Θέση της παρούσας εργασίας}
\label{sec:bg-positioning}

Η βιβλιογραφία που παρουσιάστηκε παραπάνω είναι, ως επί το πλείστον,
αρχιτεκτονική/επισκοπητική (\cite{li2025advancing, majamaa2024toward}) ή
εστιάζει σε συγκεκριμένα ζητήματα πρωτοκόλλου
(\cite{noh2024tnntn}), και βασίζεται σε αναλυτικά μοντέλα stochastic
geometry ή σε προσομοιώσεις επιπέδου σύστησης (system-level) με ιδιαίτερη
έμφαση στην κάλυψη (coverage probability) ως συνάρτηση κατωφλίου SINR
(βλ. την περίπτωση μελέτης πολυσυνδεσιμότητας έναντι μονο-συνδεσιμότητας
στο~\cite{li2025advancing}, ή τα αναλυτικά μοντέλα κάλυψης cooperative LEO
δικτύων στο~\cite{shang2023coverage}).

Η παρούσα εργασία διαφοροποιείται ως εξής: αντί για αναλυτικά μοντέλα
κάλυψης ή high-level περιγραφή αρχιτεκτονικών, αναπτύσσει μια
\emph{αναλυτική, per-link προσομοίωση} πάνω σε συγκεκριμένη γεωμετρία
(πραγματικές γεωγραφικές συντεταγμένες σταθμών βάσης, χρηστών και
δορυφόρου), με πρότυπα μοντέλα διάδοσης 3GPP TR~38.901 για το επίγειο
σκέλος, και ρητή μοντελοποίηση ενεργειακής κατανάλωσης. Αυτό επιτρέπει την
εξαγωγή συγκεκριμένων, αναπαραγώγιμων αριθμητικών KPI (όχι μόνο πιθανότητα
κάλυψης, αλλά throughput και ενέργεια ανά bit ανά χρήστη), και -- κυρίως --
την επέκταση σε χρονική διάσταση με πραγματική κίνηση του δορυφόρου, κάτι
που δεν καλύπτεται από τα παραπάνω άρθρα παρά μόνο σε επίπεδο συζήτησης
(π.χ. η αναφορά των Majamaa~\cite{majamaa2024toward} σε \emph{quasi-earth
fixed} έναντι \emph{earth-moving} beam deployments, χωρίς αριθμητική
προσομοίωση). Η προσομοίωση αυτή αποτελεί επίσης τον γεννήτορα δεδομένων
για την πρώτη προσπάθεια εφαρμογής μηχανικής μάθησης στην απόφαση επιλογής
κόμβου, θέμα που δεν θίγεται καθόλου στα~\cite{li2025advancing,
noh2024tnntn, majamaa2024toward}.

\section{Μοντέλα διάδοσης καναλιού 3GPP TR~38.901}
\label{sec:bg-channel-model}

Το πρότυπο 3GPP TR~38.901~\cite{tr38901} ορίζει στοχαστικά γεωμετρικά
μοντέλα καναλιού για συχνότητες από 0.5 έως 100~GHz, καλύπτοντας σενάρια
όπως Urban Macro (UMa), Urban Micro Street Canyon (UMi), Rural Macro (RMa)
κ.ά. Δύο στοιχεία του προτύπου είναι κεντρικά για τη μεθοδολογία της
παρούσας εργασίας (αναλυτικά στην Ενότητα~\ref{sec:meth-terrestrial}):

\begin{itemize}
  \item Η \textbf{πιθανότητα οπτικής επαφής} (Line-of-Sight - LOS) μεταξύ
    πομπού και δέκτη, η οποία στο TR~38.901 (§7.4.2) εκφράζεται ως
    συνάρτηση της οριζόντιας απόστασης $d_{2D}$ και, στο σενάριο UMa, και
    του ύψους του χρήστη, και χρησιμοποιείται για μια στοχαστική
    (Bernoulli) απόφαση LOS/NLOS ανά ζεύξη, αντί για μια σταθερή, καθολική
    παραδοχή.
  \item Το \textbf{shadow fading}, μοντελοποιημένο ως log-normal
    τυχαία μεταβλητή γύρω από τη μέση απώλεια διάδοσης, με τυπική απόκλιση
    που εξαρτάται από το σενάριο και την κατάσταση LOS/NLOS (§7.4.1).
\end{itemize}

Η στοχαστικότητα αυτών των δύο στοιχείων -- ένα διαφορετικό δείγμα LOS και
shadow fading ανά ζεύξη και ανά εκτέλεση -- είναι αυτό ακριβώς που κάνει
ουσιαστική τη στατιστική ανάλυση επαναλαμβανόμενων εκτελέσεων στην
Ενότητα~\ref{sec:results-repeated}, και αυτό που, όπως θα δειχθεί στην
Ενότητα~\ref{sec:results-temporal}, γίνεται προβληματικό όταν εφαρμόζεται
\emph{ανεξάρτητα σε κάθε χρονικό βήμα} μιας χρονικής προσομοίωσης, καθώς
δεν αντιστοιχεί σε μια φυσικά συσχετισμένη, χρονικά εξελισσόμενη
πραγματικότητα.

Για το δορυφορικό σκέλος, η παρούσα εργασία δεν χρησιμοποιεί το πλήρες
μοντέλο NTN καναλιού του TR~38.821~\cite{tr38821} (το οποίο περιλαμβάνει
επιπλέον απώλειες λόγω ατμόσφαιρας/σπινθηρισμού και κερδών κεραίας), αλλά
ένα απλοποιημένο μοντέλο απωλειών ελεύθερου χώρου (Free-Space Path Loss -
FSPL) με μάσκα ελάχιστης γωνίας ανύψωσης -- μια συνειδητή απλοποίηση που
συζητείται περαιτέρω ως κατεύθυνση μελλοντικής εργασίας στο
Κεφάλαιο~\ref{ch:discussion}.

\section{Ενεργειακή μοντελοποίηση σταθμών βάσης}
\label{sec:bg-energy-model}

Για τη μοντελοποίηση της κατανάλωσης ισχύος του επίγειου σταθμού βάσης, η
εργασία υιοθετεί το γραμμικό μοντέλο ισχύος EARTH των
Auer~κ.ά.~\cite{auer2011earth}, $P = N_{\text{TRX}}(P_0 + \Delta_P
P_{\text{out}})$ σε ενεργή λειτουργία, όπου $P_0$ είναι η σταθερή
κατανάλωση ισχύος ανεξάρτητα από το φορτίο κίνησης, $\Delta_P$ η κλίση
κατανάλωσης ως προς την εκπεμπόμενη ισχύ $P_{\text{out}}$, και
$N_{\text{TRX}}$ ο αριθμός πομποδεκτών ανά σταθμό βάσης. Το μοντέλο αυτό
παραμένει σημείο αναφοράς στη βιβλιογραφία ενεργειακής απόδοσης δικτύων,
ακριβώς επειδή αποσυνδέει τη σταθερή από τη μεταβλητή-ως-προς-φορτίο
κατανάλωση, επιτρέποντας ρεαλιστικές συγκρίσεις ενέργειας ανά bit μεταξύ
κόμβων με πολύ διαφορετική κλίμακα ισχύος εκπομπής -- όπως ακριβώς
συμβαίνει μεταξύ ενός επίγειου σταθμού βάσης (43~dBm) και ενός δορυφόρου
LEO (34~dBm EIRP payload, αλλά με πολύ μικρότερη σταθερή κατανάλωση
υποδομής) στο σενάριο της παρούσας εργασίας.

\section{Μηχανική μάθηση για διαχείριση συνδεσιμότητας}
\label{sec:bg-ml}

Η χρήση μηχανικής μάθησης για αποφάσεις επιλογής κόμβου, handover ή
κατανομής πόρων σε STINs αναφέρεται συχνά ως μελλοντική κατεύθυνση στα
άρθρα επισκόπησης (π.χ. learning-based channel prediction και precoding
στο~\cite{li2025advancing}, AI/ML-based scheduling
στο~\cite{majamaa2024toward}), χωρίς όμως να αναπτύσσεται πλήρως ως
εφαρμογή. Η παρούσα εργασία κάνει ένα πρώτο, μικρής κλίμακας βήμα προς
αυτή την κατεύθυνση: χρησιμοποιεί τυπικούς ταξινομητές (Λογιστική
Παλινδρόμηση, Random Forest, υλοποιημένους με τη βιβλιοθήκη
scikit-learn~\cite{pedregosa2011sklearn}) πάνω σε ένα σύνολο δεδομένων
που παράγεται από το ίδιο το προσομοιωμένο μοντέλο link-budget, με στόχο
να επαληθευτεί η ορθότητα του συνολικού pipeline δεδομένων πριν επενδυθεί
προσπάθεια σε πιο σύνθετους, ρεαλιστικούς στόχους πρόβλεψης (βλ.
Ενότητα~\ref{sec:results-ml} και Κεφάλαιο~\ref{ch:discussion}).
